\documentclass[luatex]{ltjsarticle}
\usepackage{luatexja}
\usepackage{luatexja-fontspec}
\usepackage{amsmath,amssymb,amsthm,mathtools}
\usepackage{tikz}
\usetikzlibrary{positioning,decorations.pathmorphing}

\title{S3\_CH: 二部グラフ構成と Lean 検証}
\author{TeX執筆: CODEX (OpenAI) \\ Lean 実装: CODEX (OpenAI)}
\date{2025年9月25日}

\begin{document}
\maketitle

\section{問題設定の整理}
課題 S3\_CH は、頂点集合を \(V = \{0,1,2,3,4,5\} \cong \mathrm{Fin}\,6\) とする単純グラフで、完全二部グラフ \(K_{3,3}\) から辺 \(\{0,3\}\) を 1 本取り除いたものが二部グラフであり、さらに辺数が 8 であることを示す問題である。二部性は頂点分割 \(L = \{0,1,2\},\ R = \{3,4,5\}\) によって証明できる。

\section{具体的構成}
完全二部グラフ \(K_{3,3}\) の辺集合は \(L\) と \(R\) の直積全体であり 9 本存在する。本問題では以下の 8 本の辺のみを残す。
\[
E = \{(0,4),(0,5),(1,3),(1,4),(1,5),(2,3),(2,4),(2,5)\}.
\]
辺 \(\{0,3\}\) を除くことで、依然として全ての辺が \(L\) から \(R\) への接続となり、同一側に属する頂点を結ぶ辺は存在しない。

\section{二部性の確認}
頂点分割 \((L,R)\) が二部性を与えることは、任意の辺が必ず \(L\) の頂点と \(R\) の頂点を結んでいることを確かめればよい。集合 \(E\) を成分ごとに確かめれば明らかであり、取り除いた \(\{0,3\}\) 以外の辺はすべて条件を満たしている。

\section{辺数の検算}
\(K_{3,3}\) には 9 本の辺が存在し、そこから 1 本を取り除いたので \(|E| = 8\) である。集合として重複がないことは単純なケース分けで確かめられる。

\section{Lean での実装概要}
Lean ファイル \texttt{S3A.lean} では、次のステップで証明が構成されている。
\begin{enumerate}
  \item \textbf{具体的辺集合の定義}: リスト \texttt{almostK33EdgesList} から \texttt{Finset} を生成し、重複がないことを決定手続き \texttt{decide} で確定する。
  \item \textbf{単純グラフ \texttt{almostK33} の構築}: `SimpleGraph` 構造体に隣接関係を与え、対称性とループが無いことを証明する。ループ無しは補助補題 \texttt{no\_loop} を通じて個別に確認する。
  \item \textbf{二部性補題 \texttt{edge\_forces\_partition}}: 任意の辺が \(L\) と \(R\) の相対位置に従うことを有限ケース分解と \texttt{decide} で示す。
  \item \textbf{主定理 \texttt{S3\_CH}}: 上記の補題を組み合わせ、二部性と辺数 8 の同時成立を与える存在定理として結論付ける。
\end{enumerate}

\section{図式による可視化}
図\ref{fig:almostK33} に、頂点分割と欠損した辺を含む構造を示す。

\begin{figure}[h]
  \centering
  \begin{tikzpicture}[scale=1.0, every node/.style={circle, draw, inner sep=1pt}]
    \node (l0) at (0,2) {0};
    \node (l1) at (0,0) {1};
    \node (l2) at (0,-2) {2};
    \node (r3) at (4,2) {3};
    \node (r4) at (4,0) {4};
    \node (r5) at (4,-2) {5};

    % remaining edges
    \foreach \u/\v in {l0/r4,l0/r5,l1/r3,l1/r4,l1/r5,l2/r3,l2/r4,l2/r5}
      \draw (\u) -- (\v);

    % dashed missing edge
    \draw[dashed, thick, red] (l0) -- (r3);

    \node[draw=none] at (-0.7,2) {$L$};
    \node[draw=none] at (4.7,2) {$R$};
    \node[draw=none] at (2,-2.7) {$\text{欠損辺 } \{0,3\}$};
  \end{tikzpicture}
  \caption{完全二部グラフ \(K_{3,3}\) から辺 \(\{0,3\}\) を除いたグラフ}
  \label{fig:almostK33}
\end{figure}

\section{まとめ}
具体的な辺集合を固定し有限検査を用いることで、二部性および辺数の主張は迅速に確認できる。Lean 上でも同様の構造を反映し、`Fin` の場合分けと `decide` に頼ることで自動化された検証が行える。

\end{document}
